\documentclass[11pt]{article}
\usepackage[utf8]{inputenc}

%----- Configuración del estilo del documento------%
\usepackage{epsfig,graphicx}
\usepackage[left=2cm,right=2cm,top=1.8cm,bottom=2.3cm]{geometry}
\usepackage{fancyhdr}
\usepackage{lastpage}
\usepackage{url}
\usepackage{longtable}
\pagestyle{fancy}
\fancyhf{}
\fancyhead[R]{\ifnum\value{page}>1 \textbf{\'Algebra L\'ineal II}\fi}
\rfoot{\textit{Página \thepage}}

%Encabezado izquierdo: logo
\lfoot{\includegraphics[height=1.0cm]{Logo_FC.png}}



%------ Paquetes matemáticos básicos --------%
\usepackage{amsmath}
\usepackage{amssymb}
\usepackage{amsthm}
\usepackage[spanish]{babel}
\usepackage{graphicx}
\usepackage{hyperref}

\usepackage{tabularx}
\usepackage{xcolor}
\usepackage[table]{xcolor}
\usepackage{colortbl}
\usepackage{array, multirow, multicol, tabularx}
\usepackage{tcolorbox}
\newtheorem{theorem}{Theorem}[section]
\newtheorem{corollary}{Corollary}[theorem]
\newtheorem{lemma}[theorem]{Lemma}

%------si-------%
\definecolor{B}{HTML}{FFFFFF}
\definecolor{G}{HTML}{5e5e5e}
\definecolor{R2}{HTML}{d53d40}
\definecolor{A2}{HTML}{034190}
\definecolor{V2}{HTML}{7faa50}

\newcommand{\PrdI}[2]{\ensuremath{\langle #1, #2 \rangle}}
\newcommand{\R}{\mathbb{R}}
\newcommand{\C}{\mathbb{C}}
\newcommand{\N}{\mathbb{N}}
\newcommand{\Z}{\mathbb{Z}}
\newcommand{\Q}{\mathbb{Q}}
\renewcommand{\theenumi}{\Roman{enumi}}
\renewcommand{\labelenumi}{{\theenumi}.}

\begin{document}
	\makeatletter
        \renewenvironment{proof}[1][\proofname]{\par
            \pushQED{\qed}%
            \normalfont \topsep6\p@\@plus6\p@\relax
            \trivlist
            \item\relax
            {\itshape
            #1\@addpunct{.}}\par\vspace{\baselineskip}\ignorespaces
            }{%
            \popQED\endtrivlist\@endpefalse
            }
        \newenvironment{Solucion}[1][Soluci\'on]{\par
            \pushQED{\qed}%
            \normalfont \topsep6\p@\@plus6\p@\relax
            \trivlist
            \item\relax
            {\itshape
            #1\@addpunct{.}}\par\vspace{\baselineskip}\ignorespaces
            }{%
            \popQED\endtrivlist\@endpefalse
            }
            
    \makeatother

%------ Encabezado -------- %

\begin{center}
    \begin{minipage}{3cm}
    	\begin{center}
    		\includegraphics[height=3.4cm]{logo_unam.png}
    	\end{center}
    \end{minipage}\hfill
    \begin{minipage}{10cm}
    	\begin{center}
    	\textbf{\large Universidad Nacional Autónoma de México}\\[0.1cm]
        \textbf{Facultad de Ciencias}\\[0.1cm]
        \textbf{\'Algebra L\'ineal II}\\[0.1cm]
         Tarea examen 5 \\[0.1cm]
         El\'ias L\'opez Rivera\\[0.1cm]
        \texttt{elias.lopezr,@ciencias.unam.mx }\\[0.1cm]
        Fecha:\,\,04/12/2025
    	\end{center}
    \end{minipage}\hfill
    \begin{minipage}{3cm}
    	\begin{center}
    		\includegraphics[height=3.4cm]{Logo_FC.png}
    	\end{center}
    \end{minipage}
\end{center}

\rule{17cm}{0.1mm}

%------ Fin de encabezado -------- %
\begin{tcolorbox}[title = \textcolor{black}{\textcolor{white}{Problema 1}}]
\textit{Sea $T$ el operador l\'ineal sobre $P_2(\R)$ definido por:\,\\
\begin{equation*}
    T(f(x))=2f(x)-f'(x)
\end{equation*}
Encuentre una base para cada subespacio propio generalizado de $T$ que consista en una uni\'on de ciclos ajenos de vectores propios generalizados.Luego determine la forma canonica de Jordan }
\end{tcolorbox}
\begin{proof}
    Tomemos la base canonica $\beta=\{1,x,x^2\}$ luego se tiene que:\,\\
    \begin{align*}
        T(1)=2(1)-0=2\,\\
        \\
        T(x)=2x-1\,\\
        \\
        T(x)=2x^2-2x\,\\
    \end{align*}
    Por tanto la matriz que representa al operador respecto a la base $\beta$:\,\\
    \begin{equation*}
       A=[T]_{\beta_1}=\begin{pmatrix}
            2&-1&0\\
            0&2&-2\\
            0&0&2
        \end{pmatrix}
    \end{equation*}
    \,\\
    Encontrando el polinomio m\'inomo caracteristico para $T$:\,\\
    \begin{equation*}
        det[A-\lambda Id]=\begin{vmatrix}
            2-\lambda&-1&0\\
            0&2-\lambda&-2\\
            0&0&2-\lambda
        \end{vmatrix}=(\lambda-2)^3
    \end{equation*}\,\\
    Luego encontrando el polinomio m\'inimo para $T$:\,\\
    \begin{align*}
        [A-2id]^2&=\begin{pmatrix}
            0&-2&0\\
            0&0&1\\
            0&0&0\\
        \end{pmatrix}
        \begin{pmatrix}
            0&-2&0\\
            0&0&1\\
            0&0&0\\
        \end{pmatrix}=\begin{pmatrix}
            0&0&2\\
            0&0&0\\
            0&0&0
        \end{pmatrix}\\
        \,\\
        [A-2id]^3&=\begin{pmatrix}
            0&0&2\\
            0&0&0\\
            0&0&0
        \end{pmatrix}
        \begin{pmatrix}
            0&-2&0\\
            0&0&1\\
            0&0&0\\
        \end{pmatrix}=\begin{pmatrix}
            0&0&0\\
            0&0&0\\
            0&0&0
        \end{pmatrix}
    \end{align*}\,\\
    Por tanto el polinomio m\'inimo es $\mu_T(\lambda)=(\lambda-2)^3$, es decir que la longitud del m\'aximo bloque de jordan coincide con la dimensi\'on del
    espacio, por tanto la forma canonica de Jordan es:\,\\
    \begin{equation*}
        J([T]_{\beta_1})=\begin{pmatrix}
            2&1&0\\
            0&2&1\\
            0&0&2
        \end{pmatrix}
    \end{equation*}
    Para encontar la base de Jordan para $T$ debemos encontrar un $3$-ciclo para ello calculamos:\,\\
    \begin{align*}
        \begin{pmatrix}
            0&-1&0\\
            0&0&-2\\
            0&0&0
        \end{pmatrix}
        \begin{pmatrix}
            a\\
            b\\
            c\\
        \end{pmatrix}&=\begin{pmatrix}
            -b\\
            -2c\\
            0
        \end{pmatrix}\implies
        b=c=0\,\\
        \,\\
        Ker(A-2Id)&=\left\{\begin{pmatrix}
            a\\
            0\\
            0
        \end{pmatrix}\,\,a\in \R\right\}\,\\
        \,\\
        \beta_2&=\left\{\begin{pmatrix}
        1\\
        0\\
        0
        \end{pmatrix}\right\}\xrightarrow{base}Ker(A-2Id)\\
        \,\\
          \begin{pmatrix}
            0&0&2\\
            0&0&0\\
            0&0&0
        \end{pmatrix}
        \begin{pmatrix}
            a\\
            b\\
            c\\
        \end{pmatrix}&=\begin{pmatrix}
            c\\
            0\\
            0
        \end{pmatrix}\implies
        c=0\,\\
        \,\\
        Ker(A-2Id)^2&=\left\{\begin{pmatrix}
            a\\
            b\\
            0
        \end{pmatrix}\,\,a,b\in \R\right\}\,\\
        \,\\
        \beta_3&=\left\{\begin{pmatrix}
        1\\
        0\\
        0
        \end{pmatrix}\begin{pmatrix}
        0\\
        1\\
        0
        \end{pmatrix}\right\}\xrightarrow{base}Ker(A-2Id)^2
    \end{align*}
    Por tanto tomamos $e_3=\begin{pmatrix}
        0\\
        0\\
        1
    \end{pmatrix}$ es claro que $e_3\notin Ker(A-2Id)^2$,$e_3\notin ker(A-2Id)$, luego tomamos el $3$-ciclo para $f=x^2$:\,\\
    \begin{align*}
        (T-2Id)(x^2)&=T(x^2)=T(x^2)-2x=-2x\\
        \,\\
        (T-2Id)^2(x^2)&=(T-2Id)(-2x)=-4x-2+4x=-2
    \end{align*}\,\\
    Por tanto la base de Jordan es $J:=\{-2,-2x,x^2\}$
\end{proof}\,\\
\begin{tcolorbox}[title = \textcolor{black}{\textcolor{white}{Problema 2}}]
\textit{Sea $T$ el operador l\'ineal sobre $M_{2\times 2}(\R)$ definido por:\,\\
\begin{equation*}
    T(A)=\begin{pmatrix}
        1&1\\
        0&1
    \end{pmatrix}A\,\,\,\,A\in M_{2\times 2}(\R)
\end{equation*}
Encuentre bases para los subespacios propios generalizados y determine una forma can\'onica de Jordan para $T$
}
\end{tcolorbox}
\begin{proof}
    Tomemos la base can\'onica de $M_{2\times 2}(\R)$:\,\\
    \begin{equation*}
        \beta_1=\left\{\begin{pmatrix}
            1&0\\
            0&0
        \end{pmatrix},\begin{pmatrix}
            0&1\\
            0&0
        \end{pmatrix},\begin{pmatrix}
            0&0\\
            1&0
        \end{pmatrix},\begin{pmatrix}
            0&0\\
            0&1
        \end{pmatrix}\right\}
    \end{equation*}\,\\
    Luego tenemos que:\,\\
    \begin{align*}
        T\begin{pmatrix}
            1&0\\
            0&0
        \end{pmatrix}&=\begin{pmatrix}
            1&1\\
            0&1
        \end{pmatrix}\begin{pmatrix}
            1&0\\
            0&0
        \end{pmatrix}=\begin{pmatrix}
            1&0\\
            0&0
        \end{pmatrix}\\
        \,\\
        T\begin{pmatrix}
            0&1\\
            0&0
        \end{pmatrix}&=\begin{pmatrix}
            1&1\\
            0&1
        \end{pmatrix}\begin{pmatrix}
            0&1\\
            0&0
        \end{pmatrix}=\begin{pmatrix}
            0&1\\
            0&0
        \end{pmatrix}\\
        \,\\
        T\begin{pmatrix}
            0&0\\
            1&0
        \end{pmatrix}&=\begin{pmatrix}
            1&1\\
            0&1
        \end{pmatrix}\begin{pmatrix}
            0&0\\
            1&0
        \end{pmatrix}=\begin{pmatrix}
            1&0\\
            1&0
        \end{pmatrix}\,\\
        \,\\
           T\begin{pmatrix}
            0&0\\
            0&1
        \end{pmatrix}&=\begin{pmatrix}
            1&1\\
            0&1
        \end{pmatrix}\begin{pmatrix}
            0&0\\
            0&1
        \end{pmatrix}=\begin{pmatrix}
            0&1\\
            0&1
        \end{pmatrix}
    \end{align*}\,\\
    Se sigue que la matriz de $T$ respecto a la base $\beta$ es:\,\\
    \begin{equation*}
      A=[T]_{\beta_1}=\begin{pmatrix}
        1&0&1&0\\
        0&1&0&1\\
        0&0&1&0\\
        0&0&0&1
        \end{pmatrix}
    \end{equation*}
    Encontrando el polinomio caracteristico para $T$:
    \begin{equation*}
       p_{T}(\lambda)= det[A-\lambda Id]=\begin{vmatrix}
        1-\lambda&0&1&0\\
        0&1-\lambda&0&1\\
        0&0&1-\lambda&0\\
        0&0&0&1-\lambda
        \end{vmatrix}=(\lambda-1)^4
    \end{equation*}\,\\
    Luego encontrando el polinomio m\'inimo para $T$:\,\\
    \begin{align*}
        [A-Id]^2&=\begin{pmatrix}
        0&0&1&0\\
        0&0&0&1\\
        0&0&0&0\\
        0&0&0&0
        \end{pmatrix}
        \begin{pmatrix}
        0&0&1&0\\
        0&0&0&1\\
        0&0&0&0\\
        0&0&0&0
        \end{pmatrix}=\begin{pmatrix}
        0&0&0&0\\
        0&0&0&0\\
        0&0&0&0\\
        0&0&0&0
        \end{pmatrix}
    \end{align*}\,\\
    Por tanto $m_T(\lambda)=(\lambda-1)^2$, es decir el ciclo de longitud m\'as grande sera un $2-ciclo$, tenemos que las posibles particiones en ciclos del operador son dos $2$-ciclos y un $2$ ciclo con dos $1$-ciclo, para determinar
    que partici\'on es la correcta determinamos la dimensi\'on del subespacio generalizado la cual nos dira la cantidad de bloques necesarios:\,\\
    \begin{align*}
        Ker(A-Id)&=\begin{pmatrix}
            0&0&1&0\\
            0&0&0&1\\
            0&0&0&0\\
            0&0&0&0\\
        \end{pmatrix}\,\begin{pmatrix}
            b_{1,1}\\
            b_{1,2}\\
            b_{2,1}\\
            b_{2,2}
        \end{pmatrix}=\begin{bmatrix}
            b_{2,1}\\
            b_{2,2}\\
            0\\
            0\\
        \end{bmatrix}\implies b_{2,1}=b_{2,2}=0\\
        \,\\
        Ker(A-id)&=\left\{\begin{pmatrix}
        b_{1,1}\\
        b_{2,2}\\
        0\\
        0
        \end{pmatrix}\,\,b_{1,1},b_{2,2}\in \R\right\}\\
        \,\\
        \beta_2:&=\left\{\begin{pmatrix}
        1\\
        0\\
        0\\
        0
        \end{pmatrix},\begin{pmatrix}
        0\\
        1\\
        0\\
        0
        \end{pmatrix}\right\}\xrightarrow{base}Ker(A-Id)
    \end{align*}\,\\
    Como el subespacio propio tiene dimensi\'on dos esto implica que se necesitan dos bloques de jordan, por tanto la partici\'on correcta es en 
    dos $2$-ciclos:\,\\
    \begin{equation*}
        J([T]_{\beta_1})=\begin{pmatrix}
            1&1&0&0\\
            0&1&0&0\\
            0&0&1&1\\
            0&0&0&1
        \end{pmatrix}
    \end{equation*}
\end{proof}
\begin{tcolorbox}[title = \textcolor{black}{\textcolor{white}{Problema 3}}]
\textit{Sea $T:\R^5\rightarrow \R^5$ un operador l\'ineal cuyo polinomio caracter\'istico es:
\begin{equation*}
    p_T(\lambda)=(\lambda-2)^3\,(\lambda+1)^2
\end{equation*}
y cuyo polinomio m\'inimo es
\begin{equation*}
    m_T(\lambda)=(\lambda-2)^2(\lambda+1)
\end{equation*}\,\\
Determine todas las posibles formas can\'onicas de Jordan para $T$}
\end{tcolorbox}
\begin{proof}
    Sabemos que para el valor propio $2$ el tamaño del mayor ciclo es $2$ pues la multiplicidad algerabica de $(\lambda-2)$ en $\mu_T$ es $2$ mientras que el tamaño de todos los ciclos para el valor propio $2$ debe sumar $3$ pues e la multiplicidad de $(\lambda-2)$ en $p_T$ es $3$, la unica partici\'on posible
    es en un $2$-ciclo y  un $1$-ciclo, de forma similar el ciclo de mayor longitud para el valor propio $-1$ es $1$ pues es la multiplicidad de $(\lambda+1)$ en $\mu_t$ es $1$, como la longitud de los ciclos para este valor propio debe sumar $2$ pues es la multiplicidad de $(\lambda+1)$ en $p_T$, la \'unica partici\'on posible es en dos $1$-ciclo,por tanto la \'unica forma de Jordan posible es:\,\\
    \begin{equation*}
        J(T)=\begin{pmatrix}
            2&1&0&0&0\\
            0&2&0&0&0\\
            0&0&2&0&0\\
            0&0&0&1&0\\
            0&0&0&0&1
        \end{pmatrix}
    \end{equation*}
\end{proof}
\begin{tcolorbox}[title = \textcolor{black}{\textcolor{white}{Problema 4}}]
\textit{Sea $T:\R^4\rightarrow \R^4$ un operador l\'ineal cuyo polinomio caracter\'istico es:
\begin{equation*}
    p_T(\lambda)=(\lambda-1)^4
\end{equation*}
y cuyo polinomio m\'inimo es
\begin{equation*}
    m_T(\lambda)=(\lambda-1)^2
\end{equation*}\,\\
Determine todas las posibles formas can\'onicas de Jordan para $T$ y todas las descomposiciones posibles de ciclos de vectores propios generalizados}
\end{tcolorbox}
\begin{proof}
    Tenemos que las posibles dimesniones para $E_{\lambda_1}$ son $\{1,2,3,4\}$, es claro que esta no puede ser $4$, pues el operador seria diagonalizable, ni $1$ ya que el bloque de mayor tamaño tiene longitud $2$ pues es la multiplicidad de $(\lambda-1)$ en $\mu_T$, por lo que si hubiera un solo bloque no se alcanzaria la dimensi\'on del espacio,
    es decir solo puede haber $2$ o $3$ bloques en total, como necesariamente tenemos un $2$-ciclo, las particiones adecuadas para poder alcanzar la dimensi\'on del espacio son, un $2$-ciclo y dos $1$-ciclos ($3$ bloques) \'o dos $2$-ciclos ($2$ bloques), en cuyo caso tenemos dos formas de Jordan posibles:\,\\
    \,\\
    \begin{equation*}
        J_1(T)=\begin{pmatrix}
            1&1&0&0\\
            0&1&0&0\\
            0&0&1&0\\
            0&0&0&1
        \end{pmatrix}
    \end{equation*}\,\\
    \begin{equation*}
        J_2(T)=\begin{pmatrix}
            1&1&0&0\\
            0&1&0&0\\
            0&0&1&1\\
            0&0&0&1
        \end{pmatrix}
    \end{equation*}  
\end{proof}
\begin{tcolorbox}[title = \textcolor{black}{\textcolor{white}{Definici\'on}}]
\textit{Un operador l\'ineal se dice nilpotente si existe un entero positvio $k$ tal que:
\begin{equation*}
    N^k=0
\end{equation*}
el menor entero que cumple estya propiedad se llama \textbf{\'indice de nilpotencia}.Los operadores nilpotentes
tienen \'unicamente el valor propio $0$ y su forma de Jordan se compone de bloques:
\begin{equation*}
    J_{0,m}=[0,e_1,e_2,\cdots,e_{m-1}]
\end{equation*}}
\end{tcolorbox}
\,\\
\begin{tcolorbox}[title = \textcolor{black}{\textcolor{white}{Problema 5}}]
\textit{Sea $N$ un operador nilpotente de \'indice de nilpotencia $m$. Demuestre que para cada entero $k\geq 1$
\begin{equation*}
    dim Nuc(N^k)-dim Nuc(N^{k-1})
\end{equation*}
es igual al n\'umero de bloques de Jordan de tama\~no al menos $k$ en la forma de Jordan de $N$.
Concluya que el operador nilpotente queda completamente determinado (hasta semejanza) por la sucesi\'on:
\begin{equation*}
    dim Nuc(N),dim Nuc(N^2),\cdots,dim Nuc(N^m)
\end{equation*}}
\end{tcolorbox}
\begin{proof}
    Consideremos que el operador $N$ tiene forma de Jordan:\,\\
    \,\\
    \begin{equation*}
        [N]_{\beta}=\begin{pmatrix}
            J_{0,k_1}&0&\cdots&0\\
            0&J_{0,k_2}&\cdots&0\\
            \vdots&&\vdots\\
            0&0&\cdots&J_{0,k_d}
        \end{pmatrix}
    \end{equation*}\,\\
Es claro que $m<dim(V)=n$.
\begin{lemma}
    Sea $n\in \Z^{+}$, sea $J_{0,n}$ el bloque de Jordan para el valor propio $0$ de tamaño $n$, tenemos que para 
    $k\in\{1,2,\cdots,n\}$ se cumple que $dim(Im(J^k_{0,n}))=n-k$, para valores de $k$ mayores a $n$ se tiene que $dim(Im(J^k_{0,n}))=0$
\end{lemma}
\begin{proof}
    Si $\{e_1,\cdots,e_n\}$ es la base canonica, es claro que $J_{0,n}e_i=e_{i-1}$ para $i\in\{2,\cdots,n\}$ y para $J_{n,0}e_1=0$, es decir, la multiplicaci\'on por $J_{0,n}$ va desplazando los elementos de la base canonica hacia a la izquierda y reemplazando por columnas de ceros, entonces si hacemos este proceso 
    suficientes veces obtendremos la matriz cero, de este modo para $k\in\{1,2,\cdots,n\}$ se obtiene lo siguiente:\,\\
    \begin{equation*}
        J_{0,n}^k=\begin{cases}
            0\,\,\text{para}\,\,k\geq i\\
            e_{i-k}\,\,\text{para}\,\,k\leq i-1
        \end{cases}
    \end{equation*}\,\\
    Asi $J_{0,n}^k$ env\'ia a los primeros $k$ vectores de la base canonica al vector cero, por tanto la imagen de $J_{0,n}^k$ est\'a generada por $\{e_{k+1},e_{k+2},\cdots,e_n\}$, luego $dim(Im(J_{0,n}^k))=n-k$.
\end{proof}
\,\\
Ahora es claro que para $r\in\N$ se cumple que:\,\\
    \,\\
    \begin{equation*}
        [N]^r_{\beta}=\begin{pmatrix}
            J_{0,k_1}^r&0&\cdots&0\\
            0&J_{0,k_2}^r&\cdots&0\\
            \vdots&&\vdots\\
            0&0&\cdots&J_{0,k_d}^r
        \end{pmatrix}
    \end{equation*}\,\\
    Ahora definimos:\,\\
    \begin{equation*}
        N_r:=\{J_{0,k_i}\,\,es\,\,un\,\,bloque\,\,de\,\,N\,\,y\,\,k_i=r\}
    \end{equation*}\,\\
    ademas de $m_r=|N_r|$, es decir $m_r$ representa \textbf{la cantidad de bloques de $N$ de tamaño $r$ dentro de la forma de Jordan de $N$ }, en particular tenemos que claramente la forma de jordan de $N$ no puede contener bloques de tamaño mayor a $dim(V)=n$ por tanto reestringimos $r\leq n$, luego tenemos que:\,\\
    \begin{equation*}
        dim(Im(N^r))=\sum_{i=1}^d\,dim(Im(J_{o,k_i}^r))=\sum_{i=r}^{n}\,m_i(i-r)
    \end{equation*}\,\\
    Luego tenemos que:\,\\
    \begin{align*}
        dim(Im(N^{r-1}))-dim(Im(N^r))&=\sum_{j=r-1}^{n}m_j(j-r+1)-\sum_{i=r}^{n}\,m_i(i-r)=\sum_{j=1}^{n-r+1}jm_{r+(j-1)}+\sum_{i=1}^{n-r}im_{r+i}\\
        \,\\
        &=m_r+\sum_{i=1}^{n-r}(i+1)\,m_{r+i}-\sum_{i=1}^{n-r}im_{r+i}=m_r+\sum_{i=1}^{n-r}(i+1-i)m_{r+i}\\
        \,\\
        &=m_r+\sum_{i=1}^{n-r}m_{r+i}=\sum_{i=r}^{n}m_{i}
    \end{align*}\,\\
    Ahora se sigue que:\,\\
    \begin{align*}
        dim(Nuc(N^r))-dim(Nuc(N^{r-1}))&=dim(V)-dim(Im(N^r))-dim(V)+dim(Im(N^{r-1}))\\
        \,\\
        &=dim(Im(N^{r-1}))-dim(Im(N^r))=\sum_{i=r}^{n}m_{i}
    \end{align*}\,\\
Ahora es f\'acil ver que $dim(Nuc(N^r))-dim(Nuc(N^{r-1}))$ es igual al n\'umero de bloques de Jordan de tama\~no al menos $r$ en la forma de Jordan de $N$, ademas tenemos que no podemos tener bloquer de tamaño mayor a $m$, pues la cantidad de bloques de tamaño $m+1$ es $dim(Nuc(N^m+1))-dim(Nuc(N^{m}))=0$ pues $N^m=N^{m+1}=0$, por tanto
todos los bloques de la forma de Jordan de $N$ se encuentran definidos por la sucesi\'on:
\begin{equation*}
    dim(Nuc(N)),dim(Nuc(N^2)),\cdots,dim(Nuc(N^m)) 
\end{equation*}
\end{proof}
\end{document}